\documentclass[aspectratio=169,10pt,xcolor={dvipsnames,table}]{beamer}

% Tema Moderno Metropolis
\usetheme{metropolis}

% Pacotes Fundamentais
\usepackage[utf8]{inputenc}
\usepackage[T1]{fontenc}
\usepackage[brazilian]{babel}
\usepackage{graphicx}
\usepackage{booktabs}
\usepackage{amsmath}
\usepackage{amssymb}
\usepackage{fontawesome5}
\usepackage{tcolorbox}
\tcbuselibrary{skins,raster}
\usepackage{bookmark}
\usepackage{tikz}

% ==============================================================================
% Paleta de Cores Institucional IFCE Campus Taua + Modern Tech
% ==============================================================================
\definecolor{IFCEGreen}{RGB}{47, 158, 65}       % Verde IFCE oficial vibrante
\definecolor{IFCEDarkGreen}{RGB}{20, 102, 46}    % Verde escuro institucional elegante
\definecolor{IFCELightGreen}{RGB}{236, 253, 243} % Fundo verde claro moderno
\definecolor{IFCERed}{RGB}{200, 16, 46}          % Vermelho IFCE oficial
\definecolor{IFCELightRed}{RGB}{254, 242, 242}   % Fundo vermelho suave
\definecolor{TechDark}{RGB}{15, 23, 42}          % Slate escuro premium para texto
\definecolor{TechGray}{RGB}{71, 85, 105}         % Cinza moderno para subtitulos
\definecolor{TechLight}{RGB}{248, 250, 252}      % Fundo clean suave
\definecolor{TechBorder}{RGB}{226, 232, 240}     % Borda sutil
\definecolor{NetBlue}{RGB}{2, 132, 199}          % Azul redes / nuvem
\definecolor{NetLightBlue}{RGB}{240, 249, 255}   % Fundo azul suave
\definecolor{CyberPurple}{RGB}{124, 58, 237}     % Roxo seguranca / tecnologia
\definecolor{CyberLightPurple}{RGB}{245, 243, 255}
\definecolor{GoldAmber}{RGB}{217, 119, 6}        % Ambar / destaque visual

% Configuracao de Cores no Metropolis Beamer
\setbeamercolor{palette primary}{fg=white, bg=IFCEDarkGreen}
\setbeamercolor{palette secondary}{fg=white, bg=IFCEGreen}
\setbeamercolor{palette tertiary}{fg=white, bg=TechDark}
\setbeamercolor{structure}{fg=IFCEDarkGreen}
\setbeamercolor{titlelike}{fg=white}
\setbeamercolor{background canvas}{bg=white}
\setbeamercolor{normal text}{fg=TechDark}
\setbeamercolor{progress bar}{fg=IFCERed, bg=IFCEGreen!30}
\setbeamercolor{alerted text}{fg=IFCERed}
\setbeamercolor{frametitle}{fg=white, bg=IFCEDarkGreen}

% Configuracao do cabecalho e rodape do Metropolis
\metroset{
  progressbar=frametitle,
  block=fill,
  numbering=fraction
}

% Logo Institucional IFCE Taua no canto superior direito de todos os slides de conteudo
\addtobeamertemplate{frametitle}{}{%
  \begin{tikzpicture}[remember picture,overlay]
    \node[anchor=north east, yshift=-0.11cm, xshift=-0.32cm] at (current page.north east) {
      \includegraphics[height=0.68cm]{logo_ifce_taua.png}
    };
  \end{tikzpicture}%
}

% ==============================================================================
% Cards e Caixas Estilizadas (TColorBox) com Visual Moderno
% ==============================================================================
\newtcolorbox{moderncard}[2][]{
  enhanced,
  colback=white,
  colframe=TechDark!80!black,
  colbacktitle=TechDark!85!black,
  coltitle=white,
  fonttitle=\bfseries\footnotesize,
  arc=2.2mm,
  boxrule=0.7pt,
  left=2.5mm, right=2.5mm, top=2mm, bottom=2mm,
  title={#2},
  #1
}

\newtcolorbox{greencard}[2][]{
  enhanced,
  colback=white,
  colframe=IFCEDarkGreen,
  colbacktitle=IFCEDarkGreen,
  coltitle=white,
  fonttitle=\bfseries\footnotesize,
  arc=2.2mm,
  boxrule=0.7pt,
  left=2.5mm, right=2.5mm, top=2mm, bottom=2mm,
  title={#2},
  #1
}

\newtcolorbox{bluecard}[2][]{
  enhanced,
  colback=white,
  colframe=NetBlue!90!black,
  colbacktitle=NetBlue!90!black,
  coltitle=white,
  fonttitle=\bfseries\footnotesize,
  arc=2.2mm,
  boxrule=0.7pt,
  left=2.5mm, right=2.5mm, top=2mm, bottom=2mm,
  title={#2},
  #1
}

\newtcolorbox{purplecard}[2][]{
  enhanced,
  colback=white,
  colframe=CyberPurple!85!black,
  colbacktitle=CyberPurple!85!black,
  coltitle=white,
  fonttitle=\bfseries\footnotesize,
  arc=2.2mm,
  boxrule=0.7pt,
  left=2.5mm, right=2.5mm, top=2mm, bottom=2mm,
  title={#2},
  #1
}

\newtcolorbox{redcard}[2][]{
  enhanced,
  colback=white,
  colframe=IFCERed,
  colbacktitle=IFCERed,
  coltitle=white,
  fonttitle=\bfseries\footnotesize,
  arc=2.2mm,
  boxrule=0.7pt,
  left=2.5mm, right=2.5mm, top=2mm, bottom=2mm,
  title={#2},
  #1
}

\newtcolorbox{statcard}[1][]{
  enhanced,
  colback=TechLight,
  colframe=TechBorder,
  arc=2.2mm,
  boxrule=0.6pt,
  left=2mm, right=2mm, top=1.8mm, bottom=1.8mm,
  halign=center,
  #1
}

\newtcolorbox{highlightbox}[3][]{
  enhanced,
  colback=#2,
  colframe=#3,
  arc=2.2mm,
  boxrule=0.8pt,
  left=2.5mm, right=2.5mm, top=2mm, bottom=2mm,
  #1
}

% ==============================================================================
% Badges Modernas Estilo Tag
% ==============================================================================
\newcommand{\pilltag}[2]{%
  \tcbox[
    on line,
    arc=1.5mm,
    colback=#1,
    colframe=#1,
    boxrule=0pt,
    left=1.5mm, right=1.5mm, top=0.6mm, bottom=0.6mm,
    boxsep=0pt
  ]{\textcolor{white}{\tiny\bfseries #2}}%
}

\newcommand{\greenbadge}[1]{\pilltag{IFCEGreen}{#1}}
\newcommand{\redbadge}[1]{\pilltag{IFCERed}{#1}}
\newcommand{\bluebadge}[1]{\pilltag{NetBlue}{#1}}
\newcommand{\purplebadge}[1]{\pilltag{CyberPurple}{#1}}

% Metadados da Apresentacao
\title[Técnico em Redes - IFCE Tauá]{\textbf{Curso Técnico Integrado em Redes de Computadores}}
\subtitle{Ensino Médio de Excelência + Formação Tecnológica Federal no Campus Tauá}
\author[IFCE Tauá]{Instituto Federal de Educação, Ciência e Tecnologia do Ceará\texorpdfstring{\\ \textbf{Campus Tauá}}{ - Campus Tauá}}
\institute[IFCE]{Especial para Estudantes do 9º Ano do Ensino Fundamental}
\date{\today}

% ==============================================================================
% Início do Documento Modular
% ==============================================================================
\begin{document}

% Secao 1: Capa, Decisao do 9º Ano e o que e o Ensino Medio Integrado
% Secao 1: Capa, Decisao do 9º Ano e o que e o Ensino Medio Integrado

% Slide 1: Capa Customizada de Alto Impacto Visual
\begin{frame}[plain]
  \begin{columns}[c]
    \begin{column}{0.38\textwidth}
      \centering
      \includegraphics[width=0.82\textwidth]{logo_ifce_taua.png}
      \vspace{0.35cm}
      
      {\scriptsize \textbf{\textcolor{IFCEDarkGreen}{MINISTÉRIO DA EDUCAÇÃO}}} \\
      \vspace{0.06cm}
      {\tiny \textbf{Instituto Federal do Ceará -- Campus Tauá}} \\
      \vspace{0.04cm}
      {\tiny \textcolor{TechGray}{Diretoria de Ensino / Coordenação de Redes}}
    \end{column}
    
    \begin{column}{0.60\textwidth}
      \vspace{-0.1cm}
      {\scriptsize \textcolor{IFCERed}{\textbf{\faGraduationCap\ ENSINO MÉDIO INTEGRADO À EDUCAÇÃO PROFISSIONAL}}}
      \vspace{0.12cm}
      
      {\Large \textbf{\textcolor{IFCEDarkGreen}{Técnico em Redes\\de Computadores}}}
      
      \vspace{0.15cm}
      {\small \textbf{Conecte Seu Futuro ao Mundo a Partir de Tauá!}}
      
      \vspace{0.25cm}
      \begin{tcolorbox}[
        enhanced,
        colback=IFCELightGreen,
        colframe=IFCEGreen!70!black,
        arc=2.2mm,
        boxrule=0.7pt,
        left=2.8mm, right=2.8mm, top=2mm, bottom=2mm
      ]
        \scriptsize
        \textbf{\faCheckCircle\ Formação Dupla:} Ensino Médio Completo + Diploma Técnico Federal \\
        \textbf{\faClock\ Duração:} 3 anos letivos no Campus Tauá (Matutino / Vespertino) \\
        \textbf{\faGift\ 100\% Gratuito:} Sem mensalidade, sem taxas e com assistência estudantil
      \end{tcolorbox}
      
      \vspace{0.15cm}
      {\tiny \textcolor{TechDark!70}{\faMapMarker\ IFCE Campus Tauá -- CE \quad | \quad Especial para Estudantes do 9º Ano}}
    \end{column}
  \end{columns}
\end{frame}

% Slide 2: A Grande Escolha do 9º Ano
\begin{frame}{A Grande Escolha do 9º Ano}{\faQuestionCircle\ Onde você quer estar nos próximos 3 anos da sua vida?}
  \begin{columns}[t]
    \begin{column}{0.48\textwidth}
      \begin{moderncard}{\faSchool\ Uma Escola Comum}
        \scriptsize
        \begin{itemize}\setlength{\itemsep}{0.18em}
          \item \textbf{Apenas o currículo padrão:} 3 anos de ensino médio tradicional em sala de aula convencional.
          \item \textbf{Sem profissão ao se formar:} Aos 18 anos, você sai sem nenhuma habilitação para o mercado.
          \item \textbf{Prática quase inexistente:} Poucos ou nenhum laboratório moderno de informática.
          \item \textbf{Poucas oportunidades extras:} Raras chances de bolsas de pesquisa ou intercâmbios.
        \end{itemize}
      \end{moderncard}
    \end{column}
    
    \begin{column}{0.48\textwidth}
      \begin{greencard}{\faRocket\ No IFCE Campus Tauá}
        \scriptsize
        \begin{itemize}\setlength{\itemsep}{0.18em}
          \item \textbf{Ensino Médio de Excelência:} Professores mestres e doutores com foco no ENEM.
          \item \textbf{Profissão Valorizada aos 18 anos:} Diploma de Técnico Federal reconhecido no Brasil todo.
          \item \textbf{Mão na Massa Todo Dia:} Laboratórios climatizados com roteadores, fibra e servidores.
          \item \textbf{Mundo de Oportunidades:} Bolsas em dinheiro, robótica, olimpíadas e intercâmbio!
        \end{itemize}
      \end{greencard}
    \end{column}
  \end{columns}
  
  \vspace{0.1cm}
  \begin{center}
    \scriptsize \textcolor{IFCERed}{\textbf{\faLightbulb\ No IFCE você não escolhe entre Médio ou Técnico: você tem os dois juntos e 100\% gratuito!}}
  \end{center}
\end{frame}

% Slide 3: O que é o Ensino Médio Integrado?
\begin{frame}{O Que É o Ensino Médio Integrado?}{\faLayerGroup\ Entenda como funciona a modalidade mais completa do Brasil}
  \begin{columns}[t]
    \begin{column}{0.48\textwidth}
      \begin{bluecard}{\faBookOpen\ Ensino Médio Completo}
        \tiny
        Você cursa todas as disciplinas fundamentais:
        \begin{itemize}\setlength{\itemsep}{0.12em}
          \item Português, Literatura e Redação Nota 1000
          \item Matemática e Raciocínio Lógico
          \item Física, Química e Biologia com laboratórios
          \item História, Geografia, Filosofia e Sociologia
          \item Língua Inglesa, Artes e Educação Física
        \end{itemize}
        \textbf{\textcolor{NetBlue}{\faCheckDouble\ Resultado:}} Base completa para o ENEM e Vestibulares!
      \end{bluecard}
    \end{column}
    
    \begin{column}{0.48\textwidth}
      \begin{purplecard}{\faNetworkWired\ Formação Técnica em Redes}
        \tiny
        Integrado na mesma rotina, você se forma em Tech:
        \begin{itemize}\setlength{\itemsep}{0.12em}
          \item Montagem e Manutenção de Hardware
          \item Redes de Computadores e Roteamento
          \item Servidores Linux e Windows na prática
          \item Cibersegurança e Defesa contra Hackers
          \item Programação e Desenvolvimento Web
        \end{itemize}
        \textbf{\textcolor{CyberPurple}{\faAward\ Resultado:}} Diploma de Técnico de Nível Médio!
      \end{purplecard}
    \end{column}
  \end{columns}
  
  \vspace{0.08cm}
  \begin{statcard}
    \tiny \textcolor{IFCEDarkGreen}{\textbf{\faGraduationCap\ Conclusão em 3 Anos:} Você termina aos 18 anos pronto para a faculdade ou para conquistar vagas disputadas em tecnologia!}
  \end{statcard}
\end{frame}


% Secao 2: O que e Redes de Computadores? (A magica por tras da internet)
% Secao 2: O que e Redes de Computadores? (A magica por tras da internet)

% Slide 4: Como a Internet Funciona?
\begin{frame}{Você Já Pensou em Como a Internet Funciona?}{\faWifi\ Da tela do seu celular para qualquer lugar do planeta em segundos}
  \begin{columns}[t]
    \begin{column}{0.48\textwidth}
      \begin{bluecard}{\faMobile\ O Que Você Vê Todo Santo Dia}
        \scriptsize
        \begin{itemize}\setlength{\itemsep}{0.18em}
          \item \textbf{Games Online sem Lag:} Como Free Fire e Roblox sincronizam dezenas de jogadores ao vivo?
          \item \textbf{Vídeos em 4K:} Como streams no YouTube e TikTok chegam sem travar no seu celular?
          \item \textbf{Mensagens Instantâneas:} Como uma foto no WhatsApp atravessa o mundo em milissegundos?
          \item \textbf{Segurança no PIX:} Como seu dinheiro viaja com criptografia e proteção antifraude?
        \end{itemize}
      \end{bluecard}
    \end{column}
    
    \begin{column}{0.48\textwidth}
      \begin{greencard}{\faProjectDiagram\ O Que Existe nos Bastidores}
        \scriptsize
        \begin{itemize}\setlength{\itemsep}{0.18em}
          \item \textbf{Espinha Dorsal Global:} Milhares de km de fibra óptica, cabos submarinos e satélites.
          \item \textbf{Roteadores Inteligentes:} Equipamentos que guiam cada bit na velocidade da luz.
          \item \textbf{Data Centers Gigantescos:} Servidores 24h por dia sustentando a Nuvem e a IA.
          \item \textbf{Sem Redes, Nada Disso Existe:} O celular vira só uma calculadora sem a rede de dados!
        \end{itemize}
      \end{greencard}
    \end{column}
  \end{columns}
  
  \vspace{0.1cm}
  \begin{center}
    \scriptsize \textcolor{IFCEDarkGreen}{\textbf{\faGlobe\ Redes de Computadores é a infraestrutura invisível que conecta a humanidade!}}
  \end{center}
\end{frame}

% Slide 5: Os 4 Superpoderes do Tecnico em Redes
\begin{frame}{O Que Faz um Técnico em Redes?}{\faUserAstronaut\ Os 4 pilares práticos da profissão que conecta o mundo}
  \begin{columns}[t]
    \begin{column}{0.24\textwidth}
      \begin{bluecard}{\faNetworkWired\ 1. Conexão}
        \tiny
        \textbf{Infraestrutura}
        \begin{itemize}\setlength{\itemsep}{0.12em}
          \item Redes Wi-Fi velozes.
          \item Cabeamento e fibra.
          \item Interliga empresas.
        \end{itemize}
      \end{bluecard}
    \end{column}
    
    \begin{column}{0.24\textwidth}
      \begin{greencard}{\faServer\ 2. Servidores}
        \tiny
        \textbf{Coração da Nuvem}
        \begin{itemize}\setlength{\itemsep}{0.12em}
          \item Linux e Windows.
          \item Gestão de usuários.
          \item Mantém serviços no ar.
        \end{itemize}
      \end{greencard}
    \end{column}
    
    \begin{column}{0.24\textwidth}
      \begin{redcard}{\faLock\ 3. Segurança}
        \tiny
        \textbf{Defesa Cibernética}
        \begin{itemize}\setlength{\itemsep}{0.12em}
          \item Bloqueia malwares.
          \item Configura firewalls.
          \item Protege senhas e dados.
        \end{itemize}
      \end{redcard}
    \end{column}
    
    \begin{column}{0.24\textwidth}
      \begin{purplecard}{\faCode\ 4. Código}
        \tiny
        \textbf{Automação \& Web}
        \begin{itemize}\setlength{\itemsep}{0.12em}
          \item Scripts em Python.
          \item Páginas Web e bancos.
          \item Soluções IoT e sensores.
        \end{itemize}
      \end{purplecard}
    \end{column}
  \end{columns}
  
  \vspace{0.1cm}
  \begin{statcard}
    \tiny \textcolor{TechDark}{\faWrench\ \textbf{Profissão Completa:} Você aprende tanto a parte física (hardware, cabos, placas) quanto a parte lógica (sistemas, códigos, segurança)!}
  \end{statcard}
\end{frame}


% Secao 3: Mercado de Trabalho e Portas para o Futuro (Emprego e Universidade)
% Secao 3: Mercado de Trabalho e Portas para o Futuro (Emprego e Universidade)

% Slide 6: O Mercado de Trabalho em Redes
\begin{frame}{Mercado de Trabalho: Uma Área Sempre Aquecida}{\faBriefcase\ Por que empresas disputam técnicos em Redes no Ceará e no Brasil?}
  \begin{columns}[t]
    \begin{column}{0.48\textwidth}
      \begin{greencard}{\faCity\ Onde Você Pode Trabalhar?}
        \tiny
        \begin{itemize}\setlength{\itemsep}{0.18em}
          \item \textbf{Provedores de Internet (ISPs):} Empresas de fibra óptica em Tauá e região sempre contratando técnicos de suporte e rede.
          \item \textbf{Empresas e Comércio Local:} Lojas, bancos, clínicas, hospitais, escolas e órgãos públicos que dependem de sistemas 24h.
          \item \textbf{Indústrias \& Agro 4.0:} Redes industriais de automação e fazendas conectadas com sensores inteligentes.
          \item \textbf{Trabalho Remoto (Home Office):} Empresas nacionais e multinacionais contratando profissionais à distância.
        \end{itemize}
      \end{greencard}
    \end{column}
    
    \begin{column}{0.48\textwidth}
      \begin{bluecard}{\faChartLine\ Cenário do Setor de Tecnologia}
        \tiny
        \begin{itemize}\setlength{\itemsep}{0.18em}
          \item \textbf{Alta Demanda por Técnicos:} Faltam milhares de profissionais com formação prática de qualidade no Brasil.
          \item \textbf{Setor Essencial:} Nenhuma empresa hoje fatura ou atende clientes sem internet e servidores protegidos.
          \item \textbf{Independência Financeira Cedo:} Comece a receber seu próprio salário aos 18 anos recém-saído do Ensino Médio.
          \item \textbf{Crescimento Acelerado:} De suporte júnior a administrador de redes ou especialista em cibersegurança.
        \end{itemize}
      \end{bluecard}
    \end{column}
  \end{columns}
  
  \vspace{0.08cm}
  \begin{statcard}
    \tiny \textcolor{IFCEDarkGreen}{\faAward\ \textbf{Registro no CFT:} Ao concluir, você tem direito ao registro no Conselho Federal dos Técnicos Industriais para assinar projetos!}
  \end{statcard}
\end{frame}

% Slide 7: Dupla Porta para o Futuro
\begin{frame}{Dupla Porta Para o Futuro: Você Escolhe o Seu Destino}{\faDirections\ Ao terminar o 3º ano no IFCE Tauá, todos os caminhos estão abertos}
  \begin{columns}[t]
    \begin{column}{0.48\textwidth}
      \begin{greencard}{\faDollarSign\ Opção 1: Mercado Imediato}
        \tiny
        \textbf{Comece a Trabalhar aos 18 Anos:}
        \begin{itemize}\setlength{\itemsep}{0.18em}
          \item Diploma de Técnico de Nível Médio Federal em mãos.
          \item 3 anos de experiência prática em laboratórios de verdade.
          \item Autonomia financeira para comprar suas coisas e ajudar em casa.
          \item Capacidade para prestar serviços autônomos ou abrir empresa de TI.
        \end{itemize}
      \end{greencard}
    \end{column}
    
    \begin{column}{0.48\textwidth}
      \begin{purplecard}{\faUniversity\ Opção 2: Faculdade de Topo}
        \tiny
        \textbf{Acesso Direto ao Ensino Superior:}
        \begin{itemize}\setlength{\itemsep}{0.18em}
          \item Preparação forte para o ENEM com professores de ponta.
          \item Ingresso em ADS no IFCE Tauá, Ciência da Computação ou Engenharias.
          \item Vantagem absurda: você já domina Linux, redes e programação!
          \item Conquista imediata de estágios remunerados na faculdade.
        \end{itemize}
      \end{purplecard}
    \end{column}
  \end{columns}
  
  \vspace{0.08cm}
  \begin{center}
    \scriptsize \textcolor{IFCERed}{\textbf{\faStar\ E você pode trabalhar E fazer faculdade ao mesmo tempo, bancando seus estudos com a própria profissão!}}
  \end{center}
\end{frame}


% Secao 4: A Jornada dos 3 Anos (Matriz Curricular ano a ano)
% Secao 4: A Jornada dos 3 Anos (Matriz Curricular ano a ano)

% Slide 8: Visao Geral da Jornada Formativa
\begin{frame}{A Jornada dos 3 Anos: Do Zero ao Especialista}{\faRoute\ Um percurso passo a passo planejado para quem nunca mexeu em servidores}
  \begin{columns}[t]
    \begin{column}{0.32\textwidth}
      \begin{greencard}{\faSeedling\ 1º Ano: Fundações}
        \tiny
        \textbf{Descobrindo o Mundo Tech}
        \begin{itemize}\setlength{\itemsep}{0.12em}
          \item Como computadores conversam.
          \item Dominando Linux e Windows.
          \item Eletricidade básica e hardware.
          \item Base forte de matérias do Médio.
        \end{itemize}
        \vspace{0.06cm}
        \centering\greenbadge{Nível: Iniciante Curioso}
      \end{greencard}
    \end{column}
    
    \begin{column}{0.32\textwidth}
      \begin{bluecard}{\faTools\ 2º Ano: Mão na Massa}
        \tiny
        \textbf{Construindo a Infraestrutura}
        \begin{itemize}\setlength{\itemsep}{0.12em}
          \item Cabeamento estruturado e racks.
          \item Configuração prática de Servidores.
          \item Lógica de Programação em Python.
          \item Desmontagem e conserto de PCs.
        \end{itemize}
        \vspace{0.06cm}
        \centering\bluebadge{Nível: Técnico Prático}
      \end{bluecard}
    \end{column}
    
    \begin{column}{0.32\textwidth}
      \begin{purplecard}{\faLock\ 3º Ano: Elite \& Nuvem}
        \tiny
        \textbf{Segurança, Wi-Fi e Web}
        \begin{itemize}\setlength{\itemsep}{0.12em}
          \item Cibersegurança e Firewalls.
          \item Redes Sem Fio e antenas de rádio.
          \item Páginas Web e Banco de Dados.
          \item Prática Profissional e portfólio.
        \end{itemize}
        \vspace{0.06cm}
        \centering\purplebadge{Nível: Especialista Pronto}
      \end{purplecard}
    \end{column}
  \end{columns}
  
  \vspace{0.08cm}
  \begin{statcard}
    \tiny \textcolor{IFCEDarkGreen}{\faLightbulb\ \textbf{Pergunta Frequente:} ``Preciso saber programar ou consertar computador para entrar?'' $\rightarrow$ \textbf{Não!} Você aprende tudo do zero com os professores!}
  \end{statcard}
\end{frame}

% Slide 9: 1º Ano
\begin{frame}{1º Ano: O Ponto de Partida no Universo Tech}{\faCompass\ Os primeiros passos na transição do Ensino Fundamental para o Médio Técnico}
  \begin{columns}[t]
    \begin{column}{0.50\textwidth}
      \begin{greencard}{\faMicrochip\ Disciplinas Técnicas do 1º Ano (280h)}
        \tiny
        \begin{itemize}\setlength{\itemsep}{0.15em}
          \item \textbf{Comunicação de Dados e Redes (120h):} Como dados viajam por fios e ondas; modelos de rede, pacotes e endereçamento IP.
          \item \textbf{Sistemas Operacionais (80h):} O que está por trás do Windows e a mágica do terminal Linux; arquivos e processos.
          \item \textbf{Eletricidade Básica (80h):} Circuitos elétricos, medições com multímetro, fontes de alimentação e segurança.
          \item \textbf{Introdução ao Curso e Orientação (40h):} Carreira em TI, ética e rotina de estudos no ensino federal.
        \end{itemize}
      \end{greencard}
    \end{column}
    
    \begin{column}{0.48\textwidth}
      \begin{moderncard}{\faGraduationCap\ Ensino Médio de Excelência (600h)}
        \tiny
        \begin{itemize}\setlength{\itemsep}{0.15em}
          \item \textbf{Língua Portuguesa \& Redação (80h):} Domínio da escrita e interpretação crítica de textos.
          \item \textbf{Matemática (120h):} Álgebra, geometria e raciocínio lógico.
          \item \textbf{Ciências da Natureza (120h):} Física (40h), Química (40h) e Biologia (40h) com práticas experimentais.
          \item \textbf{Ciências Humanas (160h):} História, Geografia, Filosofia e Sociologia.
          \item \textbf{Linguagens:} Espanhol (40h), Artes (40h) e Ed. Física (40h).
        \end{itemize}
      \end{moderncard}
    \end{column}
  \end{columns}
  
  \vspace{0.06cm}
  \begin{center}
    \tiny \textcolor{TechGray}{\faInfoCircle\ Carga Horária Total do 1º Ano: 960h letivas para uma adaptação acolhedora e estimulante.}
  \end{center}
\end{frame}

% Slide 10: 2º Ano
\begin{frame}{2º Ano: Subindo de Nível (Infra, Servidores \& Código)}{\faCubes\ O ano da transformação: muito laboratório, hardware e as primeiras linhas de código}
  \begin{columns}[t]
    \begin{column}{0.50\textwidth}
      \begin{bluecard}{\faServer\ Disciplinas Técnicas do 2º Ano (360h)}
        \tiny
        \begin{itemize}\setlength{\itemsep}{0.15em}
          \item \textbf{Infraestrutura e Projetos de Redes (80h):} Cabeamento estruturado, conexões de fibra, patch panels e racks.
          \item \textbf{Servidores (120h):} Instalação e administração de servidores DNS, DHCP, Web e arquivos em ambientes de rede.
          \item \textbf{Montagem e Manutenção (80h):} Diagnóstico completo de hardware, montagem de PCs e manutenção preventiva/corretiva.
          \item \textbf{Lógica de Programação (80h):} Algoritmos, variáveis, repetições e criação de programas em Python para automação.
        \end{itemize}
      \end{bluecard}
    \end{column}
    
    \begin{column}{0.48\textwidth}
      \begin{greencard}{\faLightbulb\ Inovação \& Ensino Médio (680h)}
        \tiny
        \begin{itemize}\setlength{\itemsep}{0.15em}
          \item \textbf{Projeto Integrador (40h):} Criação de soluções em equipe para problemas da região dos Inhamuns.
          \item \textbf{Empreendedorismo (40h):} Modelos de negócio digitais, startups e precificação de serviços de TI.
          \item \textbf{Língua Inglesa (40h):} O idioma global da documentação tech.
          \item \textbf{Base BNCC (600h):} Português, Matemática e Ciências com foco na preparação para o ENEM.
        \end{itemize}
      \end{greencard}
    \end{column}
  \end{columns}
  
  \vspace{0.06cm}
  \begin{center}
    \tiny \textcolor{NetBlue}{\textbf{\faLaptopCode\ No 2º ano você já é capaz de montar seu próprio computador e configurar uma rede corporativa inteira!}}
  \end{center}
\end{frame}

% Slide 11: 3º Ano
\begin{frame}{3º Ano: A Elite da Conectividade (Segurança, Wi-Fi e Web)}{\faUserShield\ O ciclo de conclusão: tecnologias avançadas, defesa cibernética e preparação para o futuro}
  \begin{columns}[t]
    \begin{column}{0.50\textwidth}
      \begin{purplecard}{\faLock\ Disciplinas Técnicas do 3º Ano (320h)}
        \tiny
        \begin{itemize}\setlength{\itemsep}{0.15em}
          \item \textbf{Gerência e Segurança de Redes (80h):} Firewalls corporativos, criptografia, detecção de invasões e segurança digital.
          \item \textbf{Redes Sem Fio (80h):} Wi-Fi corporativo de alta performance, propagação de sinal, antenas e proteção wireless.
          \item \textbf{Introdução ao Desenvolvimento Web (80h):} Páginas modernas e responsivas em HTML, CSS e JavaScript.
          \item \textbf{Banco de Dados (80h):} Modelagem e consultas SQL para armazenar dados com velocidade e segurança.
        \end{itemize}
      \end{purplecard}
    \end{column}
    
    \begin{column}{0.48\textwidth}
      \begin{redcard}{\faMedal\ Reta Final: Prática \& ENEM (680h)}
        \tiny
        \begin{itemize}\setlength{\itemsep}{0.15em}
          \item \textbf{Prática Profissional Supervisionada (40h):} Vivência real em projetos com supervisão docente.
          \item \textbf{Foco Total no ENEM:} Língua Portuguesa e Redação (120h), Matemática (80h) e Biologia (80h).
          \item \textbf{Disciplinas Optativas:} Oficinas de redação e preparatórios para Olimpíadas de Física e Astronomia.
          \item \textbf{Formatura Dupla:} Conclusão do Médio com registro profissional de Técnico Federal!
        \end{itemize}
      \end{redcard}
    \end{column}
  \end{columns}
  
  \vspace{0.06cm}
  \begin{center}
    \tiny \textcolor{IFCEDarkGreen}{\textbf{\faTrophy\ Resultado Final: 3.000 horas de capacitação integral que transformam você em um profissional tech disputado pelo mercado!}}
  \end{center}
\end{frame}


% Secao 5: Mao na Massa (Laboratorios de Ponta, Racks e Projetos)
% Secao 5: Mao na Massa (Laboratorios de Ponta, Racks e Projetos)

% Slide 12: Laboratorios de Ponta
\begin{frame}{Mão na Massa: Laboratórios de Padrão Industrial}{\faMicrochip\ Nada de ficar só olhando pro quadro: você aprende fazendo com equipamentos reais}
  \begin{columns}[t]
    \begin{column}{0.48\textwidth}
      \begin{greencard}{\faNetworkWired\ Lab de Redes \& Telecomunicações}
        \tiny
        \begin{itemize}\setlength{\itemsep}{0.18em}
          \item \textbf{Racks Industriais Reais:} Os mesmos gabinetes e estruturas usados por grandes provedores e data centers.
          \item \textbf{Equipamentos de Ponta:} Switches gerenciáveis, roteadores Gigabit Dual-Band e centrais VoIP.
          \item \textbf{Bancada de Ferramental:} Alicates de crimpagem, testadores de cabos de rede e fusão de fibra óptica.
          \item \textbf{Simulação \& Prática:} Aprenda a capturar pacotes e auditar o tráfego com Wireshark e Cisco Packet Tracer.
        \end{itemize}
      \end{greencard}
    \end{column}
    
    \begin{column}{0.48\textwidth}
      \begin{bluecard}{\faDesktop\ Laboratórios de Informática 1 e 2}
        \tiny
        \begin{itemize}\setlength{\itemsep}{0.18em}
          \item \textbf{Mais de 60 Computadores:} Máquinas modernas e rápidas para cada aluno praticar de forma individual.
          \item \textbf{Conforto Térmico Total:} Ambientes 100\% climatizados com ar-condicionado e estações ergonômicas.
          \item \textbf{Internet Ultrarrápida Federal:} Conexão de fibra em alta velocidade para downloads e pesquisas.
          \item \textbf{Acessibilidade Completa:} Postos adaptados para Pessoas com Deficiência (PcD).
        \end{itemize}
      \end{bluecard}
    \end{column}
  \end{columns}
  
  \vspace{0.08cm}
  \begin{statcard}
    \tiny \textcolor{IFCEDarkGreen}{\faCheckCircle\ \textbf{Diferencial Único:} Na escola comum o computador só serve para digitar trabalhos; no IFCE você aprende como ele funciona por dentro e como fazê-lo conversar com o mundo!}
  \end{statcard}
\end{frame}

% Slide 13: Projetos Reais, Robotica e TECTEL
\begin{frame}{Projetos Reais, Robótica e Eventos Científicos}{\faFlask\ Desenvolva ideias inovadoras e apresente soluções para a sua comunidade}
  \begin{columns}[t]
    \begin{column}{0.48\textwidth}
      \begin{purplecard}{\faRobot\ Robótica \& Internet das Coisas (IoT)}
        \tiny
        \begin{itemize}\setlength{\itemsep}{0.18em}
          \item \textbf{Dispositivos Inteligentes:} Conecte placas Arduino e ESP32 com sensores à rede Wi-Fi.
          \item \textbf{Automação Residencial e Rural:} Sistemas de irrigação automática, sensores de temperatura e controle remoto por celular.
          \item \textbf{Projetos Integradores:} Equipes de alunos criando protótipos funcionais para desafios reais de Tauá.
        \end{itemize}
      \end{purplecard}
    \end{column}
    
    \begin{column}{0.48\textwidth}
      \begin{greencard}{\faCalendar\ Eventos \& Feiras de Ciência}
        \tiny
        \begin{itemize}\setlength{\itemsep}{0.18em}
          \item \textbf{TECTEL (Semana da Tecnologia):} O maior evento tech da região, com oficinas práticas, palestras e minicursos.
          \item \textbf{Hackathons \& Maratonas:} Desafios estimulantes para criar soluções de programação em equipe.
          \item \textbf{Viagens Técnicas:} Visitas guiadas a provedores de internet, emissoras e empresas de tecnologia.
        \end{itemize}
      \end{greencard}
    \end{column}
  \end{columns}
  
  \vspace{0.08cm}
  \begin{center}
    \scriptsize \textcolor{IFCEDarkGreen}{\textbf{\faTrophy\ Destaque Regional: Os estudantes do IFCE Tauá frequentemente conquistam medalhas em feiras científicas estaduais e nacionais!}}
  \end{center}
\end{frame}


% Secao 6: Vida no Campus (Bolsas em Dinheiro, Intercambio Exterior e Esportes)
% Secao 6: Vida no Campus (Bolsas em Dinheiro, Intercambio Exterior e Esportes)

% Slide 14: Bolsas em Dinheiro (PIBIC-Jr)
\begin{frame}{Bolsas em Dinheiro: Ganhe Para Estudar e Pesquisar!}{\faMoneyBillWave\ Quantas escolas de Ensino Médio pagam bolsa mensal para você estudar?}
  \begin{columns}[t]
    \begin{column}{0.48\textwidth}
      \begin{greencard}{\faFlask\ O Que É o PIBIC-Jr?}
        \tiny
        \begin{itemize}\setlength{\itemsep}{0.18em}
          \item \textbf{Iniciação Científica Júnior:} Programa de pesquisa com bolsas financiadas pelo CNPq e pelo IFCE.
          \item \textbf{Bolsa Mensal em Dinheiro:} Pagamento direto na conta bancária do estudante durante o projeto.
          \item \textbf{Mentoria de Doutores:} Você pesquisa lado a lado com professores mestres e doutores do campus.
          \item \textbf{Temas de Fronteira:} Redes de sensores, segurança digital, IA e internet das coisas no semiárido.
        \end{itemize}
      \end{greencard}
    \end{column}
    
    \begin{column}{0.48\textwidth}
      \begin{bluecard}{\faHandsHelping\ Bolsas de Extensão e Monitoria}
        \tiny
        \begin{itemize}\setlength{\itemsep}{0.18em}
          \item \textbf{Bolsas de Extensão:} Projetos que levam a informática e oficinas digitais para a comunidade de Tauá.
          \item \textbf{Monitoria de Disciplinas:} Ajude colegas em sala e laboratório recebendo auxílio financeiro mensal.
          \item \textbf{Independência Financeira:} Comece a ter seu próprio dinheiro para comprar seus livros e equipamentos!
        \end{itemize}
      \end{bluecard}
    \end{column}
  \end{columns}
  
  \vspace{0.08cm}
  \begin{statcard}
    \tiny \textcolor{IFCEDarkGreen}{\textbf{\faPiggyBank\ Transformação Real:} É a chance de juntar seu dinheiro, apoiar sua família e valorizar seu currículo antes dos 18 anos!}
  \end{statcard}
\end{frame}

% Slide 15: Intercambio Internacional
\begin{frame}{Intercâmbio Internacional: O Mundo a Partir de Tauá!}{\faPlaneDeparture\ Estude um semestre no exterior com tudo 100\% pago pelo IFCE}
  \begin{columns}[t]
    \begin{column}{0.52\textwidth}
      \begin{bluecard}{\faGlobeAmericas\ Experiência Global Sem Sair do Bolso}
        \tiny
        \begin{itemize}\setlength{\itemsep}{0.18em}
          \item \textbf{Tradição Comprovada:} Desde 2016 o IFCE Campus Tauá seleciona discentes para estudar fora do país!
          \item \textbf{Destinos de Prestígio:} Estudantes de Tauá já foram para Portugal e Canadá com bolsa integral!
          \item \textbf{Custeio Completo:} Passagens aéreas de ida e volta, seguro-saúde e recursos para hospedagem e refeições.
          \item \textbf{Auxílio Pré-Embarque:} Recurso institucional para emissão de passaporte e vistos.
        \end{itemize}
      \end{bluecard}
    \end{column}
    
    \begin{column}{0.44\textwidth}
      \begin{tcolorbox}[
        enhanced,
        colback=IFCELightRed,
        colframe=IFCERed,
        arc=2mm,
        boxrule=0.7pt,
        left=2mm, right=2mm, top=1.8mm, bottom=1.8mm
      ]
        \tiny
        \textbf{\textcolor{IFCERed}{\faPassport\ Oportunidade Inédita:}} \\
        \vspace{0.06cm}
        Quantas escolas convencionais oferecem a chance real de fazer intercâmbio no exterior com bolsa durante o Ensino Médio?
      \end{tcolorbox}
      \vspace{0.08cm}
      \begin{statcard}
        \tiny \textcolor{TechDark}{\faStar\ \textbf{Currículo de Peso:} Uma vivência no exterior aos 17 anos abre portas no mercado tech internacional!}
      \end{statcard}
    \end{column}
  \end{columns}
\end{frame}

% Slide 16: Esportes, Cultura e Vida Estudantil
\begin{frame}{Esportes, Cultura e Amizades Para a Vida Toda}{\faFutbol\ Porque a vida no IFCE é completa, dinâmica e inesquecível}
  \begin{columns}[t]
    \begin{column}{0.48\textwidth}
      \begin{greencard}{\faRunning\ Quadra Poliesportiva \& JIFCE}
        \tiny
        \begin{itemize}\setlength{\itemsep}{0.18em}
          \item \textbf{Quadra Coberta Moderna:} Futsal, vôlei misto, handebol e basquete nos intervalos e após as aulas.
          \item \textbf{Clube de Xadrez:} Prática de xadrez consagrada com torneios e medalhas no campus.
          \item \textbf{Jogos dos Institutos (JIFCE):} Viaje pelo Ceará e pelo Brasil representando Tauá nos jogos esportivos!
        \end{itemize}
      \end{greencard}
    \end{column}
    
    \begin{column}{0.48\textwidth}
      \begin{purplecard}{\faBookReader\ Biblioteca \& Convivência}
        \tiny
        \begin{itemize}\setlength{\itemsep}{0.18em}
          \item \textbf{Biblioteca Climatizada:} Acervo de livros físicos, milhares de títulos digitais e cabines de estudo.
          \item \textbf{Áreas Verdes:} Espaços arborizados com bancos para conversar, relaxar e fazer novas amizades.
          \item \textbf{Protagonismo Estudantil:} Grêmio Estudantil ativo, gincanas culturais e acolhimento contínuo.
        \end{itemize}
      \end{purplecard}
    \end{column}
  \end{columns}
  
  \vspace{0.08cm}
  \begin{center}
    \scriptsize \textcolor{IFCEDarkGreen}{\textbf{\faSmileBeam\ A Melhor Época da Vida: No IFCE você faz amizades sinceras com quem sonha grande como você!}}
  \end{center}
\end{frame}


% Secao 7: Assistencia Estudantil e Acolhimento (Alimentacao, Saude e Permanencia)
% Secao 7: Assistencia Estudantil e Acolhimento (Alimentacao, Saude e Permanencia)

% Slide 17: Assistencia Estudantil
\begin{frame}{Assistência Estudantil: O IFCE Apoia a Sua Família}{\faHands\ Nenhum estudante fica para trás por dificuldades financeiras}
  \begin{columns}[t]
    \begin{column}{0.48\textwidth}
      \begin{bluecard}{\faHandHoldingUsd\ Auxílios Financeiros Diretos}
        \tiny
        \begin{itemize}\setlength{\itemsep}{0.18em}
          \item \textbf{Auxílio Transporte:} Suporte financeiro para o deslocamento de quem mora longe ou nos distritos.
          \item \textbf{Alimentação Escolar \& Auxílio:} Refeições gratuitas com cardápio planejado por nutricionista.
          \item \textbf{Auxílio Moradia:} Para estudantes que vêm de municípios vizinhos da região dos Inhamuns.
          \item \textbf{Auxílio Óculos \& Didático:} Recursos para exames de vista, óculos de grau e material escolar.
        \end{itemize}
      \end{bluecard}
    \end{column}
    
    \begin{column}{0.48\textwidth}
      \begin{greencard}{\faCheckCircle\ Compromisso Social Federal}
        \tiny
        \begin{itemize}\setlength{\itemsep}{0.18em}
          \item \textbf{100\% Gratuito:} Sem mensalidade, taxa de matrícula ou cobranças escondidas.
          \item \textbf{Permanência com Dignidade:} Editais semestrais transparentes com foco em quem mais precisa.
          \item \textbf{Tranquilidade Para os Pais:} O jovem estuda em um ambiente seguro, com refeição e acolhimento.
          \item \textbf{Material e Conectividade:} Apoio contínuo para ninguém perder atividades por falta de recurso.
        \end{itemize}
      \end{greencard}
    \end{column}
  \end{columns}
  
  \vspace{0.08cm}
  \begin{statcard}
    \tiny \textcolor{IFCEDarkGreen}{\faCheckDouble\ \textbf{Objetivo Institucional:} Garantir que você entre, aprenda com conforto e se forme com um futuro transformado!}
  \end{statcard}
\end{frame}

% Slide 18: Saude e Acolhimento no Campus
\begin{frame}{Saúde e Acolhimento: Cuidado Completo Com Você}{\faHeartbeat\ Uma equipe multiprofissional gratuita dentro do próprio campus}
  \begin{columns}[t]
    \begin{column}{0.48\textwidth}
      \begin{purplecard}{\faUserMd\ Equipe de Saúde no Campus}
        \tiny
        \begin{itemize}\setlength{\itemsep}{0.18em}
          \item \textbf{Psicologia Escolar:} Acolhimento emocional, combate à ansiedade e adaptação aos estudos.
          \item \textbf{Assistência Social:} Acompanhamento socioeconômico e orientação contínua.
          \item \textbf{Enfermagem \& Consultório Odontológico:} Atendimento básico de saúde e odontologia no instituto.
          \item \textbf{Nutricionista:} Alimentação escolar saudável e balanceada para o dia a dia.
        \end{itemize}
      \end{purplecard}
    \end{column}
    
    \begin{column}{0.48\textwidth}
      \begin{greencard}{\faUniversalAccess\ Apoio Pedagógico \& Inclusão}
        \tiny
        \begin{itemize}\setlength{\itemsep}{0.18em}
          \item \textbf{Coordenação Pedagógica:} Orientação individual de rotina e métodos de estudo.
          \item \textbf{NAPNE (Núcleo de Inclusão):} Acompanhamento de Pessoas com Deficiência e adaptação didática.
          \item \textbf{Monitorias \& Reforço:} Horários de plantão para tirar dúvidas nas matérias mais puxadas.
          \item \textbf{Ambiente Acolhedor:} Respeito mútuo, combate ao preconceito e escuta ativa.
        \end{itemize}
      \end{greencard}
    \end{column}
  \end{columns}
  
  \vspace{0.08cm}
  \begin{center}
    \scriptsize \textcolor{IFCEDarkGreen}{\textbf{\faHandHoldingHeart\ Aqui você é acolhido como ser humano integral, com apoio, incentivo e respeito!}}
  \end{center}
\end{frame}


% Secao 8: Mitos vs. Verdades, Como Ingressar (35 Vagas) e Encerramento
% Secao 8: Mitos vs. Verdades, Como Ingressar (35 Vagas) e Encerramento

% Slide 19: Mitos vs. Verdades
\begin{frame}{Mitos vs. Verdades Sobre Estudar no IFCE Tauá}{\faCheckSquare\ Vamos tirar todas as dúvidas que passam pela cabeça de quem está no 9º ano!}
  \begin{columns}[t]
    \begin{column}{0.48\textwidth}
      \begin{redcard}{\faTimesCircle\ Mitos Que Você Já Deve Ter Ouvido}
        \tiny
        \begin{itemize}\setlength{\itemsep}{0.18em}
          \item \textbf{Mito: ``É só para quem é gênio da matemática!''} \\
          \textcolor{TechGray}{Mentira! Ensinamos tudo com calma a partir do zero.}
          \item \textbf{Mito: ``Tem que pagar mensalidade ou taxa.''} \\
          \textcolor{TechGray}{Zero reais! É 100\% gratuito mantido pelo Governo Federal.}
          \item \textbf{Mito: ``Não sobra tempo para se divertir.''} \\
          \textcolor{TechGray}{Mentira! O campus tem esportes, música, jogos e lazer.}
          \item \textbf{Mito: ``Redes é só puxar fio na parede.''} \\
          \textcolor{TechGray}{Redes é cibersegurança, servidores, nuvem e programação!}
        \end{itemize}
      \end{redcard}
    \end{column}
    
    \begin{column}{0.48\textwidth}
      \begin{greencard}{\faCheckCircle\ As Verdades Que Você Precisa Saber}
        \tiny
        \begin{itemize}\setlength{\itemsep}{0.18em}
          \item \textbf{Verdade: Todo mundo é acolhido!} \\
          \textcolor{TechDark}{Professores atenciosos e monitores prontos para ajudar.}
          \item \textbf{Verdade: Você sai na frente no mercado!} \\
          \textcolor{TechDark}{Aos 18 anos você já tem diploma técnico federal.}
          \item \textbf{Verdade: Bolsas e intercâmbio existem de fato!} \\
          \textcolor{TechDark}{Colegas de Tauá já foram para Portugal e Canadá!}
          \item \textbf{Verdade: A infraestrutura é de outro nível!} \\
          \textcolor{TechDark}{Laboratórios climatizados, quadra coberta e refeitório.}
        \end{itemize}
      \end{greencard}
    \end{column}
  \end{columns}
  
  \vspace{0.08cm}
  \begin{statcard}
    \tiny \textcolor{IFCEDarkGreen}{\faHeart\ \textbf{Pode confiar:} Quem estuda no IFCE Tauá se apaixona pelo campus e nunca mais esquece!}
  \end{statcard}
\end{frame}

% Slide 20: Como Ingressar (Processo Seletivo)
\begin{frame}{Como Conquistar Sua Vaga: 35 Vagas Anuais}{\faDoorOpen\ O passo a passo simples para se tornar um estudante federal}
  \begin{columns}[t]
    \begin{column}{0.50\textwidth}
      \begin{greencard}{\faGraduationCap\ Processo Seletivo Regular do IFCE}
        \tiny
        \begin{itemize}\setlength{\itemsep}{0.18em}
          \item \textbf{Inscrições Online:} Feitas no portal oficial (\textbf{qselecao.ifce.edu.br}).
          \item \textbf{Critério de Seleção:} Geralmente baseado no histórico escolar do Fundamental (ou prova seletiva conforme edital).
          \item \textbf{Isenção de Taxa:} Procedimento simples para estudantes de escola pública e famílias de baixa renda.
          \item \textbf{Entrada Anual:} Turma exclusiva de 35 novos estudantes por ano!
        \end{itemize}
      \end{greencard}
      \vspace{0.06cm}
      \begin{center}
        \tiny \textcolor{IFCEDarkGreen}{\textbf{\faBullhorn\ Fique atento ao edital no 2º semestre com ajuda da sua escola!}}
      \end{center}
    \end{column}
    
    \begin{column}{0.48\textwidth}
      \begin{statcard}
        \LARGE \textcolor{IFCEDarkGreen}{\textbf{35}} \\
        \scriptsize \textbf{Vagas Anuais para Redes} \\
        \tiny \textcolor{IFCERed}{\faUsers\ Formação de excelência acompanhada de perto}
      \end{statcard}
      \vspace{0.06cm}
      \begin{bluecard}{\faBalanceScale\ Sistema de Cotas da Lei Federal}
        \tiny
        \begin{itemize}\setlength{\itemsep}{0.12em}
          \item \textbf{50\% ou mais das vagas} para alunos que fizeram o Fundamental em \textbf{Escola Pública}!
          \item Subcotas por renda, Pretos, Pardos, Indígenas e PcD.
          \item Vagas de Ampla Concorrência abertas a todos.
        \end{itemize}
      \end{bluecard}
    \end{column}
  \end{columns}
\end{frame}

% Slide 21: Encerramento de Alto Impacto
\begin{frame}[plain]
  \begin{columns}[c]
    \begin{column}{0.38\textwidth}
      \centering
      \includegraphics[width=0.82\textwidth]{logo_ifce_taua.png}
      \vspace{0.35cm}
      
      {\tiny \textbf{Rua Antônio Teixeira Benevides, 01 -- Colibris}} \\
      \vspace{0.04cm}
      {\tiny CEP: 63660-000 \quad | \quad Tauá -- CE} \\
      \vspace{0.08cm}
      {\scriptsize \textcolor{IFCEDarkGreen}{\textbf{Venha Conhecer o Campus!}}}
    \end{column}
    
    \begin{column}{0.60\textwidth}
      {\Large \textbf{\textcolor{IFCEDarkGreen}{O Seu Futuro Tech Começa Agora!}}}
      
      \vspace{0.12cm}
      {\small \textbf{Venha fazer parte da família IFCE Campus Tauá.}}
      
      \vspace{0.20cm}
      \begin{moderncard}{\faInfoCircle\ Canais Oficiais de Contato \& Redes}
        \tiny
        \begin{itemize}\setlength{\itemsep}{0.18em}
          \item \textbf{\faInstagram\ Instagram:} \textcolor{IFCERed}{\textbf{@ifcetaua}} (fique por dentro de todos os editais!)
          \item \textbf{\faGlobe\ Portal Oficial:} \textcolor{NetBlue}{\textbf{www.ifce.edu.br/taua}}
          \item \textbf{\faPhone\ Telefone / Recepção:} (88) 2134-1065
          \item \textbf{\faEnvelope\ E-mail da Coordenação:} redes.taua@ifce.edu.br
        \end{itemize}
      \end{moderncard}
      
      \vspace{0.20cm}
      \centering
      \textbf{\normalsize \textcolor{IFCERed}{\faComments\ Dúvidas? Perguntem à Vontade!}} \\
      \vspace{0.06cm}
      {\tiny \textcolor{TechDark!80}{Os professores e alunos do IFCE estão aqui para conversar com vocês!}}
    \end{column}
  \end{columns}
\end{frame}


\end{document}
